% Generated draft: requires mathematical and source review.
\documentclass[11pt]{article}
\usepackage[margin=25mm]{geometry}
\usepackage{fontspec}
\setmainfont{Latin Modern Roman}
\usepackage{amsmath,amssymb,mathtools,unicode-math}
\setmathfont{Latin Modern Math}
\usepackage{enumitem,needspace,xurl,hyperref}
\setlength{\parindent}{0pt}
\setlength{\parskip}{5pt}
\setlength{\emergencystretch}{3em}
\newcommand{\N}{\mathbb{N}}
\newcommand{\Z}{\mathbb{Z}}
\newcommand{\Q}{\mathbb{Q}}
\newcommand{\R}{\mathbb{R}}
\newcommand{\C}{\mathbb{C}}
\newcommand{\F}{\mathbb{F}}
\newcommand{\A}{\mathbb{A}}
\newcommand{\PP}{\mathbb{P}}
\newcommand{\E}{\mathbb{E}}
\newcommand{\eps}{\varepsilon}
\newcommand{\dd}{\,\mathrm{d}}
\begin{document}
\section*{UnsolvedMath ID: 10400225}
\textbf{Source problem code:} catalogue entry
\subsection{Definitions and mathematical statement}
Roberts asks for the generating function for \(q\)-spin-network evaluations. The source does not specify the network class, label conventions, evaluation normalization, or variables.

For a fixed finite graph \(\Gamma\) with edge set \(E\), a customary candidate interpretation is as follows. Let \(\mathcal C(\Gamma)\subseteq\Z_{\geq0}^{E}\) be the admissible edge colorings for the chosen quantum theory, and let \(\operatorname{ev}_q(\Gamma,\mathbf c)\) be its spin-network evaluation in a specified coefficient ring \(A\). With one formal variable \(z_e\) per edge, set
\[G_\Gamma(q;\mathbf z)=\sum_{\mathbf c\in\mathcal C(\Gamma)}\operatorname{ev}_q(\Gamma,\mathbf c)\prod_{e\in E}z_e^{c_e}.\]
Once the input conventions are fixed this is a formal power series; the desired result would give a useful explicit description or formula for it beyond restating this coefficient definition. The excerpt does not assert rationality, a product form or a differential equation, and none is added here. Nor does it say that this particular multivariate coloring series, as opposed to a one-variable weighting or a series summing over graphs, is intended.

The label admissibility, choice of quantum group, treatment of vertices, graph embedding/framing and normalization of evaluations must be supplied. Whether \(q\) is an indeterminate, a complex parameter or a root of unity affects the coefficient ring and admissible labels; no such choice is made silently. Any analytic rather than formal interpretation would also need a convergence domain.
\subsection{Short English statement}
Find an explicit generating function packaging quantum spin-network evaluations. A natural coloring-series formulation is given conditionally because the source does not identify the graphs, labels, normalization or formal variables.
\subsection{Sources}
\begin{itemize}
\item Local UnsolvedMath registry, permanent ID \(10400225\), exact source code \texttt{catalogue entry}.
\item J. Roberts, attribution in Ohtsuki, \emph{Problems on invariants of knots and 3-manifolds} (2002), Problem 12.20, PDF page 168: \url{https://msp.org/gtm/2002/04/gtm-2002-04-024s.pdf}. Recorded access 2026-07-29, pdf-font-structure extraction.
\end{itemize}
\end{document}
